\chapter{Secure Channel Protocol and Protected APDU Encoding}
\label{app:scp-data-objects}

This appendix sits next to Appendix~\ref{app:apdu-over-t0}: the T=0 appendix
describes how bytes are moved, while this appendix describes what changes when
those bytes are protected by a GlobalPlatform Secure Channel Protocol.

The main protocol sequence is described in Chapter~3. This appendix gives the
longer byte-level details that would make the chapter too dense: BER-TLV
objects used by SCP11, SCP03 KDF material, and the Amendment~D protected-APDU
framing used after channel establishment.

\section{What a Secure Channel Changes}

Without a Secure Channel, the command data received by a selected application
or by the kernel management dispatcher is the data sent by the host. The status
word returned by the command handler is the status word sent back to the host.

With an active Secure Channel, several additional rules apply:
\begin{itemize}
\item the command CLA carries the GlobalPlatform secure-messaging indication;
\item command data may be encrypted before it reaches the command handler;
\item command data and selected command header bytes are authenticated by
      \texttt{C-MAC};
\item response data may be encrypted before it leaves the secure element;
\item response data and the final status word may be authenticated by
      \texttt{R-MAC};
\item the active Security Domain instance and logical channel define which
      secure-channel session is used;
\item management commands that mutate the registry are authorized through the
      active Security Domain, not merely by the selected application.
\end{itemize}

The important architectural consequence is that protected APDUs are not a new
application ABI. They are a wire-level representation. The secure-channel layer
unwraps them into the same clear APDU model described in
Appendix~\ref{app:apdu-over-t0}, then wraps the response back into the protected
wire representation.

\section{Secure Channel Establishment Objects}

Secure Channel establishment APDUs are ordinary APDUs whose data fields carry
profile-specific cryptographic objects. SCP03 uses fixed response fields around
the host and card challenges. SCP11 uses BER-TLV Control Reference Templates
and certificate/public-key Data Objects.

\subsection{SCP03}

SCP03 is the symmetric AES-based Secure Channel. It is established with
\texttt{INITIALIZE UPDATE} followed by \texttt{EXTERNAL AUTHENTICATE}.

The APDU-level establishment commands are:
\begin{itemize}
\item \texttt{INITIALIZE UPDATE}:
      \texttt{CLA=80}, \texttt{INS=50}, \texttt{P1=key version},
      \texttt{P2=00}. The command carries the host challenge and
      the response carries the card challenge, key identifiers and card
      cryptogram material needed by the host-side oracle.
\item \texttt{EXTERNAL AUTHENTICATE}:
      \texttt{CLA=84}, \texttt{INS=82}, \texttt{P1=security level},
      \texttt{P2=00}. The command carries the host cryptogram followed by the
      transmitted C-MAC. If verification succeeds, the SCP03 session becomes
      authenticated at the requested security level, and the full C-MAC becomes
      the initial secure-messaging MAC chain.
\item \texttt{PUT KEY}: a Security Domain management command used to install
      or rotate key material used by later Secure Channel establishment.
\end{itemize}

The S8 and S16 profiles differ in challenge, cryptogram and transmitted MAC
lengths. The secure-channel layer therefore asks the active Security Domain
for the current MAC length before splitting the trailing C-MAC from the
protected APDU data field.

\subsection{SCP03 Cryptographic Material and KDF}

The static key material stored in a Security Domain for SCP03 is a triple of
AES keys \gpref{scp03}{§6.1}:
\begin{itemize}
\item \texttt{Key-ENC}, used to derive the session encryption key;
\item \texttt{Key-MAC}, used to derive the session MAC keys;
\item \texttt{Key-DEK}, used to protect sensitive key data during key loading.
\end{itemize}

The specification allows AES key sizes of 128, 192, or 256 bits for this key
set \gpref{scp03}{§6.1}. The SCP03 session derivation is deterministic: the
static keys, host challenge and card challenge are combined through the
SCP03-defined KDF to produce:
\begin{itemize}
\item \texttt{S-ENC}, used for command and response confidentiality when
      enabled;
\item \texttt{S-MAC}, used for command authentication and MAC chaining;
\item \texttt{S-RMAC}, used for response authentication when R-MAC is enabled.
\end{itemize}

SCP03 defines a counter-mode KDF using CMAC as the pseudorandom function, with
a fixed input format carrying a derivation constant, a length field and a
context \gpref{scp03}{§4.1.5}. This maps to NIST SP 800-108 for the KDF
\cite{nist800108} and NIST SP 800-38B for CMAC \cite{nist80038b}. The
secure-messaging encryption path relies on AES block-cipher modes defined by
GlobalPlatform with reference to NIST block-cipher mode recommendations
\cite{nist80038a}.

A useful naming distinction is that \texttt{C-MAC} and \texttt{R-MAC} are
GlobalPlatform protocol values carried on the wire, while CMAC is the
underlying NIST authentication primitive used to compute them. The same
distinction applies to cryptograms: they are protocol values produced from the
session KDF and CMAC machinery, not a separate asymmetric authentication
primitive.

\subsection{SCP11x: SCP11a, SCP11b and SCP11c}

SCP11 is the asymmetric GlobalPlatform secure-channel family. It uses ECC key
agreement to derive the symmetric secure-messaging keys that are then used in
the same APDU protection position as SCP03.

GlobalPlatform defines three SCP11 variants:
\begin{itemize}
\item \textbf{SCP11a}: mutual authentication between the OCE and the card.
      The OCE certificate is first submitted with
      \texttt{PERFORM SECURITY OPERATION}; the channel is then opened with
      \texttt{MUTUAL AUTHENTICATE}. Both static and ephemeral card-side
      contributions are used.
\item \textbf{SCP11b}: card authentication to the OCE only. OCE
      authentication is provided outside SCP11, for example by an application
      policy. The channel is opened with \texttt{INTERNAL AUTHENTICATE}, not
      \texttt{MUTUAL AUTHENTICATE}.
\item \textbf{SCP11c}: mutual authentication using static card-side
      authentication material. It is suitable for offline card-management
      scripts authorized by \texttt{CERT.OCE.ECKA}.
\end{itemize}

The SCP identifier data object uses tag \texttt{90}. Its value is two bytes:
the first byte is \texttt{11}; the second byte encodes the SCP11 profile and
options. In the profile bits, \texttt{00} identifies SCP11b, \texttt{01}
identifies SCP11a, and \texttt{11} identifies SCP11c. Bit~3 controls whether
identity material is included in the KDF: SCP11a/SCP11b use HostID, SIN and
SDIN, whereas SCP11c uses HostID and Card Group ID.

\subsubsection{Common SCP11 APDU Objects}

All SCP11 establishment APDUs use a BER-TLV Control Reference Template with
tag \texttt{A6}. The key-agreement CRT contains:
\begin{itemize}
\item \texttt{90 02 11 pp}: SCP identifier and parameters, where
      \texttt{pp} selects SCP11a/SCP11b/SCP11c and the identity-in-KDF option;
\item \texttt{95 01 kuq}: key usage qualifier and requested secure-messaging
      level. Common values include \texttt{34} for C-MAC/R-MAC and
      \texttt{3C} for C-MAC/C-ENC/R-MAC/R-ENC; \texttt{1C} and \texttt{74}
      are SCP11c-only values;
\item \texttt{80 01 88}: AES key type;
\item \texttt{81 01 LL}: AES key length in bytes, for example
      \texttt{10} for AES-128;
\item \texttt{84 Lh <HostID>}: optional HostID, present exactly when the
      profile parameter requests identity material in the KDF;
\item \texttt{5F49 Lq <ePK.OCE.ECKA>}: the OCE ephemeral public key.
\end{itemize}

For P-256, \texttt{ePK.OCE.ECKA} is normally the uncompressed SEC1 point:
\texttt{04 || X || Y}, i.e. 65 bytes. In short APDU notation, a minimal
AES-128 CRT with HostID omitted therefore has the following shape:

\begin{verbatim}
A6 50
   90 02 11 pp
   95 01 kuq
   80 01 88
   81 01 10
   5F 49 41 04 <64 bytes of OCE ephemeral point>
\end{verbatim}

If HostID is included, the \texttt{A6} length increases by
\texttt{2 + len(HostID)} and the CRT contains:

\begin{verbatim}
84 Lh <HostID>
\end{verbatim}

For \texttt{MUTUAL AUTHENTICATE} and \texttt{INTERNAL AUTHENTICATE},
\texttt{P1.b1--b7} select the version of \texttt{SK.SD.ECKA} and
\texttt{P2.b1--b7} select its key identifier.  Bit~8 is reserved in both
parameters and must be zero; the key identifier in \texttt{P2} is non-zero.
The selected key belongs to the addressed Security Domain.  A card must reject
an unsupported selector rather than silently substituting a default ECKA key.

The identity fields are not interchangeable APDU inputs.  Tag \texttt{84}
supplies only \texttt{HostID}, and its presence must agree exactly with bit~3
of \texttt{pp}.  When that bit is set, the Security Domain supplies its own
\texttt{SIN}/\texttt{SDIN} values for SCP11a/SCP11b.  For SCP11c it supplies
the Card Group ID obtained from tag \texttt{5F20} of
\texttt{CERT.SD.ECKA}.  The OCE cannot replace those card-owned values by
adding another CRT data object.

The successful response to the SCP11 key-agreement APDU returns:

\begin{verbatim}
5F 49 Lq <card public key material>
86 10 <16-byte receipt>
90 00
\end{verbatim}

For SCP11a and SCP11b, \texttt{5F49} contains the card ephemeral public key.
For SCP11c, \texttt{5F49} contains the card static public key used in the
receipt calculation.

\paragraph{PSO selectors and chaining.}
For the SCP11 \texttt{PERFORM SECURITY OPERATION} certificate command,
\texttt{P1.b1--b7} select the CA public-key version and \texttt{P2.b1--b7}
select the CA public-key identifier.  The two high bits have independent
meanings:
\begin{itemize}
\item \texttt{P1.b8 = 1} announces another command block belonging to the same
      certificate; \texttt{P1.b8 = 0} marks its final command block;
\item \texttt{P2.b8 = 1} announces another certificate in the certificate
      chain; \texttt{P2.b8 = 0} marks the final OCE certificate.
\end{itemize}
All command blocks of one certificate retain the same CA selector and are
concatenated in order before certificate parsing and signature verification.
Changing the selector, overflowing the announced implementation bound, or
interrupting the sequence invalidates that in-progress command chain; no
partially assembled certificate may reach the Security Domain.  A profile that
does not support intermediate certificate chains rejects
\texttt{P2.b8 = 1} with \texttt{6A86}.

\subsubsection{SCP11a Exchange}

SCP11a first authenticates the OCE key material to the Security Domain, then
opens the channel with \texttt{MUTUAL AUTHENTICATE}. A representative
single-certificate exchange is:

\begin{verbatim}
1. GET DATA: retrieve CERT.SD.ECKA
   C-APDU: 80 CA BF 21 06 A6 04 83 02 kid kvn 00
   R-APDU: BF 21 L ... 7F 21 L <CERT.SD.ECKA> ... 90 00

2. PERFORM SECURITY OPERATION: submit CERT.OCE.ECKA
   C-APDU: 80 2A kvn kid Lc 7F 21 L <CERT.OCE.ECKA> 00
   R-APDU: 90 00

3. MUTUAL AUTHENTICATE: open SCP11a
   C-APDU: 80 82 sd-kvn sd-kid Lc
           A6 L
              90 02 11 01
              95 01 3C
              80 01 88
              81 01 10
              5F 49 41 04 <64-byte ePK.OCE.ECKA>
           00
   R-APDU: 5F 49 41 04 <64-byte ePK.SD.ECKA>
           86 10 <receipt>
           90 00
\end{verbatim}

If the identity-in-KDF option is enabled, the profile byte is
\texttt{05} rather than \texttt{01}, and the CRT must also contain
\texttt{84 Lh <HostID>}. The KDF SharedInfo then concatenates the key usage,
key type, key length, \texttt{HostID-LV}, \texttt{SIN-LV} and
\texttt{SDIN-LV}. SCP11a combines two ECDH contributions: one using the
authenticated static OCE public key and the card static private key, and one
using both ephemeral keys. The KDF secret is ordered
\texttt{ShSee || ShSss}. The resulting \texttt{KeyData} is assigned, in
order, to the receipt key, \texttt{S-ENC}, \texttt{S-MAC},
\texttt{S-RMAC}, and \texttt{S-DEK}. The receipt covers the exact
\texttt{A6} CRT and OCE \texttt{5F49} objects from the command followed by
the card \texttt{5F49} object returned in the response.

Freshness and replay discipline for SCP11a follows directly from the ephemeral
exchange and the session rules: a new \texttt{MUTUAL AUTHENTICATE} terminates
any previous secure-channel session on the logical channel; stale receipts,
stale MAC chains and replayed protected commands must be rejected by the
secure-messaging state.

\subsubsection{SCP11b Exchange}

SCP11b does not use \texttt{PERFORM SECURITY OPERATION} as part of the secure
channel initiation. The card authenticates itself to the OCE; the OCE must be
authorized by another mechanism before the resulting channel is accepted for
management policy. The SCP11b APDU sequence is therefore:

\begin{verbatim}
1. GET DATA: retrieve CERT.SD.ECKA
   C-APDU: 80 CA BF 21 06 A6 04 83 02 kid kvn 00
   R-APDU: BF 21 L ... 7F 21 L <CERT.SD.ECKA> ... 90 00

2. Application-specific OCE authorization
   Example: VERIFY PIN, proprietary policy, or another selected application
   contract. This step is outside SCP11 itself.

3. INTERNAL AUTHENTICATE: open SCP11b
   C-APDU: 80 88 sd-kvn sd-kid Lc
           A6 L
              90 02 11 00
              95 01 3C
              80 01 88
              81 01 10
              5F 49 41 04 <64-byte ePK.OCE.ECKA>
           00
   R-APDU: 5F 49 41 04 <64-byte ePK.SD.ECKA>
           86 10 <receipt>
           90 00
\end{verbatim}

With identity material included in the KDF, the profile byte becomes
\texttt{04} and the CRT includes \texttt{84 Lh <HostID>}. SCP11b SharedInfo
uses the same identity shape as SCP11a: key usage, key type, key length,
\texttt{HostID-LV}, \texttt{SIN-LV} and \texttt{SDIN-LV}. The shared secret
concatenates the ephemeral/ephemeral contribution with the contribution based
on the OCE ephemeral public key and the card static private key.

The key policy difference is authorization, not wrapping. SCP11b can prove the
card to the OCE and derive secure-messaging keys, but the kernel must not treat
that as OCE authentication unless a separate Security Domain or application
policy has established it.

\subsubsection{SCP11c Exchange}

SCP11c first authenticates the OCE certificate with PSO. The subsequent
mutual-authentication APDU derives the channel using the OCE static key, the
OCE ephemeral key and the card static key.

The Security Domain computes the following KDF input:

\begin{verbatim}
ShSes = ECDH(SK.SD.ECKA, ePK.OCE.ECKA)
ShSss = ECDH(SK.SD.ECKA, PK.OCE.ECKA)
ShS   = ShSes || ShSss
\end{verbatim}

SCP11c therefore generates no card-side ephemeral key pair: the public key
returned in \texttt{5F49} is the static \texttt{PK.SD.ECKA} corresponding to
\texttt{SK.SD.ECKA}.

A representative AES-128 full secure-messaging exchange is:

\begin{verbatim}
1. Optional GET DATA before offline script creation
   C-APDU: 80 CA BF 21 06 A6 04 83 02 kid kvn 00
   R-APDU: BF 21 L ... 7F 21 L <CERT.SD.ECKA> ... 90 00

2. PERFORM SECURITY OPERATION: submit CERT.OCE.ECKA
   C-APDU: 80 2A kvn kid Lc 7F 21 L <CERT.OCE.ECKA> 00
   R-APDU: 90 00

3. MUTUAL AUTHENTICATE: open SCP11c
   C-APDU: 80 82 sd-kvn sd-kid Lc
           A6 L
              90 02 11 03
              95 01 3C
              80 01 88
              81 01 10
              5F 49 41 04 <64-byte ePK.OCE.ECKA>
           00
   R-APDU: 5F 49 41 04 <64-byte PK.SD.ECKA>
           86 10 <receipt>
           90 00
\end{verbatim}

If identity material is included in the KDF, the profile byte becomes
\texttt{07}, the CRT includes \texttt{84 Lh <HostID>}, and SharedInfo uses
\texttt{HostID-LV} plus \texttt{Card Group ID-LV}. The Card Group ID is the
length/value form of tag \texttt{5F20} from \texttt{CERT.SD.ECKA}.

SCP11c has two replay-sensitive properties:
\begin{itemize}
\item the \texttt{MUTUAL AUTHENTICATE} receipt binds the CRT, the OCE
      ephemeral key and the card public key material returned in \texttt{5F49};
\item once the channel is active, protected APDUs use the same Amendment~D
      payload-and-MAC framing as SCP03, so stale command MAC chains, stale receipts and replayed
      protected commands must fail before command dispatch.
\end{itemize}

\section{Protected APDU Encoding under Amendment D}

SCP03 Amendment~D does not encode secure messaging with the ISO~7816
\texttt{DO87}, \texttt{DO97}, \texttt{DO8E}, and \texttt{DO99} objects.  The
protected command data field is instead the command payload followed directly
by the transmitted C-MAC:

\begin{verbatim}
[clear command data | encrypted padded command data] || C-MAC
\end{verbatim}

The APDU header carries the secure-messaging class bits.  An expected response
length remains the APDU \texttt{Le}; it is not moved into a data object.  The
C-MAC calculation binds the modified command header, the protected command
data length, the protected command data, and the preceding MAC chain.  The
receiver separates the trailing profile-dependent MAC length, verifies the
C-MAC, and only then decrypts and removes padding when command encryption is
active.

When response protection is active, the response has the corresponding raw
form:

\begin{verbatim}
[clear response data | encrypted padded response data] || R-MAC || SW1 || SW2
\end{verbatim}

The status word remains the actual trailing APDU status word.  It is included
in the R-MAC calculation with the protected response data; it is not replaced
by a \texttt{DO99}.  Therefore alteration of either the response data or the
status word is detected.

The same external framing is used by the SCP11 support in oXiDe SE after its
distinct asymmetric establishment and session-key derivation.  The selected
Security Domain owns the profile state, key material, chaining values, and
security level, while the common secure-channel layer owns the split,
authentication, decryption, encryption, and dispatch boundary.  Command
handlers and selected Rustlets receive only the logical clear APDU and return a
logical clear response; neither depends on the external secure-messaging
framing.
